\chapter
[\CO{[THE]} \emph{Spezielle Patientengruppe:} Thermische Schädigungen]
{\CC{[THE]} \emph{Spezielle Patientengruppe:} Thermische Schädigungen}

\ChapterIntro
{%
~~~
}
{%
\begin{description}
  \item[Maintainer:] Sebastian Gabriel
  \item[Autoren:] Diverse
  \item[Reviewer:] Standard-Reviewprozess
  \item[Version:] {\DocVersion} ({\DocVersionInfo})
  \item[SHA1:] {\DocGitCommit}
\end{description}

{\TxTeilkapitelAASS}
}

\section{Hitzeerkrankungen}

Der menschliche Körper hält seine Kerntemperatur  im
Normalfall zwischen 36,5 und 37,5{\,\textcelsius}  konstant.
Übersteigt  die Wärmezufuhr die effektive  Wärmeabgabe,
resultiert ein Hitzeschaden: \T{Hitzekollaps} ,
\T{Hitzeerschöpfung} , \T{Hitzschlag}
 (bzw. bei  direkter Sonneneinstrahlung auf den Kopf ein
 Sonnen\-stich).

\subsection{Systemische Hitzeschäden}

\subsubsection{Hitzekollaps}
\label{topic:hitzekollaps}

\Layout{0}{\TxBeschreibung{} und \TxSymptome{}}{%
Bei einem Wärmestau, bei dem die zugeführte Wärme
größer ist, als die abgegebene, kann es zu einem so
genannten Hitzekollaps kommen. 

Die \EE{Blutgefäße} in der Haut \EE{erweitern} sich und das
Blut versackt, der \EE{Blutdruck fällt ab} und die
Durchblutung des Gehirns nimmt kurzfristig ab. Die Folge ist
eine \EE{kurze Ohnmacht}, den Patient wird "`schwarz vor
Augen"'. Lagert man ihn flach, oder eventuell mit erhöhten
Beinen, und kühlt ihn, bzw. bringt ihn an einen kühlen Ort,
so bessert sich sein Zustand recht rasch.
}{
\begin{itemize}
  \item Wärmestau
  \item Erweiterung der Blutgefäße in der Haut (maximal)
  \item RR\,${\downarrow}$ , Durchblutung des Gehirns
  kurzzeitig ${\downarrow}$
  \item kurze Ohnmacht, "`schwarz vor Augen"'
\end{itemize}
}{}{}{}




% \MassnahmenNG{Spezielle Maßnahmen}{Hitzekollaps}{}{
% 
% \begin{itemize}
%   \item Flachlagerung: ${\rightarrow}$ Gehirndurchblutung
%   wieder ${\uparrow}$, bewusstseinsklar
%   \item nach Sekunden gebessert, "`regelt sich von
%   selbst"'
%   \item Patient sitzen lassen, Flüssigkeit
% \end{itemize}
% }


\subsubsection{Hitzeerschöpfung}
\label{topic:hitzeerschoepfung}

\Layout{0}{\TxBeschreibung}
{
Bei der Hitzeerschöpfung ist die \EE{Wärmeaufnahme
dauerhaft größer als die Wärmeabgabe}. Der Körper versucht
dies durch starkes Schwitzen zu kompensieren, dadurch kommt 
es zu einem \EE{starken Wasser- und Elektrolytverlust}. Der
\EE{Blutdruck fällt}, es kommt zu einer \EE{Tachykardie}, die
Körpertemperatur bleibt jedoch noch normal. 
Hier ist es wichtig den Patienten zu
kühlen.

\vspace{0.5cm}

}{
\begin{itemize}
  \item Wärmeaufnahme {\textgreater}  Wärmeabgabe
  \item Starkes Schwitzen, Wasser- und  Elektrolytverlust
  \item Starkes Durstgefühl
  \item RR ${\downarrow}$, Tachykardie
  \item Körpertemperatur normal
\end{itemize}
}{}{}{}




% \MassnahmenNG{Spezielle Maßnahmen}{Hitzeerschöpfung}{}{
% 
% \begin{itemize}
%   \item Kühlung, Schatten
%   \item Flüssigkeitszufuhr, wenn bewusstseinsklar
%   \item Neurocheck inkl. BZ-Messung
% \end{itemize}
% }



\subsubsection{Hitzschlag}
\label{topic:hitzschlag}

\Layout{0}{\TxBeschreibung{} und \TxSymptome{}}{
Beim Hitzschlag ist nun auch die \EE{Regulation der Körpertemperatur gestört},
es sind auch \E{sehr hohe Werte} (bis 43\,\textcelsius{})
möglich. Dies bedeutet \EE{akute Lebensgefahr}. 

Der Patient hat eine zunächst trockene, rote und weiße,
später grau-bläuliche Haut, der Kopf ist oft hochrot.
Weiters kommt es zu \EE{neurologischen Symptomen}, wie massiven
Kopfschmerzen, Schwindel, Erbrechen und
Bewusstseinsstörungen; bis hin zur Bewusstlosigkeit. Der Patient 
kann Schockzeichen zeigen, evtl. auch eine schnelle, flache
Atmung.
}{
\begin{itemize}
  \item Kopfschmerz, Schwindel, Erbrechen
  \item Haut zunächst trocken, rot und heiß, später grau-blau, Kopf
  hochrot
  \item Stark  erhöhte  Körpertemperatur (bis 
  43{\,\textcelsius} möglich, \EE{Lebensgefahr!})
  \item Bewusstseinsstörung, Bewusstlosigkeit
  \item Schockzeichen
  \item schnelle, flache Atmung
\end{itemize}
}{}{}{}

% \MassnahmenNG{Spezielle Maßnahmen}{Hitzschlag}{}{
% 
% \begin{itemize}
%   \item \TxMassVit
%   \item Beengende Kleidungsstücke lockern
%   \item Flachlagerung an kühlem Ort (Schatten), Beine und Kopf hochlagern. 
%   \item Kühlung von \E{außen} (Fächer, mit Eiswürfeln abreiben etc.)
%   \item Kühlung von \E{innen}: für \EE{ansprechbare} Patienten viel kühle
%   (alkoholfreie) Flüssigkeit, evtl. Elektrolyt-Getränke
%   \item Neurocheck inkl. BZ-Messung
% \end{itemize}
% }


\subsubsection{Maßnahmen bei systemischen Hitzeschäden}

\MassnahmenWide{Spezielle Maßnahmen}{Hitzekollaps, Hitzeerschöpfung, Hitzschlag}{}{
\FormatTable

\begin{tabularx}{\linewidth}{XXX}
\TabHeading{Hitzekollaps}
&
\TabHeading{Hitzerschöpfung}
&
\TabHeading{Hitzschlag}
\\\addlinespace\midrule\addlinespace
% \begin{minipage}{\linewidth}
\begin{itemize}
  \item Flachlagerung: ${\rightarrow}$ Gehirndurchblutung
  wieder ${\uparrow}$, bewusstseinsklar
  \item nach Sekunden gebessert, "`regelt sich von
  selbst"'
  \item Patient sitzen lassen, Flüssigkeit
\end{itemize}
% \end{minipage}
&
% \begin{minipage}{\linewidth}
\begin{itemize}
  \item Kühlung, Schatten
  \item Flüssigkeitszufuhr, wenn bewusstseinsklar
  \item Neurocheck inkl. BZ-Messung
\end{itemize}
% \end{minipage}
&
% \begin{minipage}{\linewidth}
\begin{itemize}
  \item \TxMassVit
  \item Beengende Kleidungsstücke lockern
  \item Flachlagerung an kühlem Ort (Schatten), Beine und Kopf hochlagern. 
  \item Kühlung von \E{außen} (Fächer, mit Eiswürfeln abreiben etc.)
  \item Kühlung von \E{innen}: für \EE{ansprechbare} Patienten viel kühle
  (alkoholfreie) Flüssigkeit, evtl. Elektrolyt-Getränke
  \item Neurocheck inkl. BZ-Messung
\end{itemize}
% \end{minipage}
\\\addlinespace\bottomrule
\end{tabularx}
}

\clearpage % TODO temporäres clearpage
\subsection{Sonnenstich}
\label{topic:sonnenstich}

\Definition{Hitzeschaden durch Hirnerwärmung infolge von
direkter Sonneneinstrahlung auf den \textbf{ungeschützten
Kopf}. }

\Layout{0}{\TxBeschreibung{} und \TxSymptome{}}{
Der Sonnenstich ist eine Sonderform des Hitzeschadens: 
er entsteht durch \EE{direkte Sonneneinstrahlung} auf den
ungeschützten Kopf. Dadurch steigt die Temperatur im Kopf,
es kommt zu einer Reizung der Hirnhäute. 

Die Folge sind Symptome wie bei
einer Hirnhautentzündung (Meningitis) beziehungsweise andere
\EE{neurologische Symptome}, wie z.\,B. Kopfschmerz, sind die
Folge. Die \EE{Gesichts- und Kopfhaut ist hochrot und heiß}, der 
restliche Körper bleibt dagegen eher kühl. Der Patient
wirkt abgeschlagen und klagt über heftige 
Kopfschmerzen, Schwindel und Übelkeit, auch Unruhe, 
Verwirrtheitszustände und Nackensteifigkeit sind zu
beobachten. In schweren Fällen kann es sogar bis zu 
Krampfanfällen und Bewusstlosigkeit kommen.

\vspace{1.2cm}

}{
\begin{itemize}
  \item Direkte Sonneneinstrahlung auf \textbf{ungeschützten
  Kopf} ${\rightarrow}$ Reizung der Hirnhäute  
  \item Symptome ähnlich einer
  Hirnhautentzündung (Nackensteifigkeit) 
  \item Evtl. kombiniert mit anderen Hitzeschäden
  \item Gesichts- \& Kopfhaut hochrot, heiß, restlicher Körper kühl
  \item Abgeschlagenheit, heftiger Kopfschmerz, Schwindel
  \item Unruhe, Verwirrtheitszustände, Brechreiz
  \item Schwere Fälle: Krampfanfälle, Bewusstlosigkeit
\end{itemize}
}{}{}{}





\MassnahmenNG{Spezielle Maßnahmen}{Sonnenstich}{}{

\begin{itemize}
  \item Vitale Bedrohung einschätzen
  \TxBeiVit \TxMassVit
  \item Flachlagerung an kühlem Ort (Schatten), Kopf hochlagern
  \item Beengende Kleidungsstücke lockern
  \item Kühlung von außen (Fächer, mit Eiswürfeln abreiben, kalte
  Umschläge auf die Stirn etc.)
  \item Kühlung von innen: für \EE{ansprechbare} Patienten kühle
  (alkoholfreie) Flüssigkeit, evtl. Elektrolyt-Getränke
%   \item Schocklagerung u. -bekämpfung. % REVIEW 2010-05-19
  \item Neurocheck inkl. BZ-Messung
\end{itemize}
}



\clearpage % TODO temp clearpage
\section{Unterkühlung}
\label{topic:unterkuehlung}

\Definition{Bei einer Körperkerntemperatur unter
36\,\textcelsius{} spricht man von Unterkühlung.}

\begin{wrapfigure}{r}{4cm}
\includegraphics[width=4cm]{./images/teaser/huette_winter_mariazell.jpg}
\end{wrapfigure}
Ähnlich
wie beim Schock kommt es auch  bei der Unterkühlung zur
Zentralisation  des Blutes, um die  Körperkerntemperatur
aufrecht erhalten  zu können. Je niedriger diese 
Kerntemperatur jedoch sinkt, desto  schwerwiegender sind die
Folgen!

Es besteht die Gefahr, dass kaltes Blut aus der Peripherie
zum Körper hin fließt und sich mit dem noch relativ warmen
zentralisierten Blut vermischt. Es kommt dann zu einer
weiteren Unterkühlung des Körperstammes, im schlimmsstem Fall
zu
einem \T{Bergungstod} (Blutfluss in Folge von Bewegungen während des
Bergungsversuches).

\paragraph{Phasen}

\begin{itemize}
  \item \T{Kampfphase}:  36,5\,--\,34\,{\textcelsius},
  Kältezittern, Herzfrequenz steigt ({\textgreater}\,100\,\nicefrac{x}{min}), schnelle
  Atmung, Schmerzen. 
  \item \T{Erschöpfungsphase}:
  34\,--\,30\,{\textcelsius}, Bewusstseinstrübung,
  Kältestarre, Herzfrequenz sinkt wieder ({\textless}\,60\,\nicefrac{x}{min}), zu langsame Atmung,
  keine Schmerzen mehr. 
  \item \T{Kältenarkose}: 30\,--\,27\,{\textcelsius},
  Bewusstlosigkeit, unregelmäßiger Herzschlag, Atempausen. 
  \item {\textless}\,27\,{\textcelsius} ${\rightarrow}$ klinisch tot. 
\end{itemize}


\MassnahmenNG{Spezielle Maßnahmen}{Schwere Unterkühlung
(\textless\,34\,\textcelsius)}{}{

\begin{itemize}
  \item \TxMassVit
  \item Den Patient so wenig als möglich (nur soviel wie unbedingt nötig)
  bewegen (drohender Bergungstod).
  \item Schutz vor weiterer Abkühlung, \Ca{keine \underbar{aktive}
  Erwärmung!}
  \begin{itemize}
    \item Woll- und Aludecken! Mütze!
  \end{itemize}
\end{itemize}
}



Bei Unterkühlung gilt:

\Fazit{"`Nobody is dead until warm and dead!"' 
Schwer unterkühlte Patienten können auch noch nach Stunden (!) ohne
Kreislauftätigkeit erfolgreich reanimiert werden!}





